\newcommand{\ReviewProjectionInterpretation}{Table~\ref{tab_review_projection_only} evaluates the six DISTS scores independently of the point-cloud branch. Pooled SROCC ranges from 0.663 on LS-PCQA to 0.927 on APSIPA. This provides a standalone projection baseline on the same grouped folds as the fused model, rather than reusing fused predictions. Unlike the historical APSIPA IQA screening, this comparison uses the final extraction and pooled OOF metrics.}
\newcommand{\ReviewFamilyInterpretation}{Table~\ref{tab_review_pca_family} compares feature representations with the same Extra Trees pipeline and shared folds. The SROCC difference (PointPCA+RS minus PointPCA+) ranges from -0.006 to 0.001. PointPCA++ has the strongest SROCC on APSIPA, ICIP2023, and SJTU-PCQA among these controlled methods, while PointPCA$^3$ is slightly higher on UnB\_PC. PointPCA+ has lower RMSE on ICIP2023 than PointPCA++. Thus, the fusion does not dominate every method on every metric. Original PointPCA is evaluated only on APSIPA using its native 46-feature vector. These comparisons isolate representations under a common regressor and do not reproduce each historical method\textquotesingle s native regression procedure.}
\newcommand{\ReviewAgreementInterpretation}{Figure~\ref{fig_review_feature_agreement} and Table~\ref{tab_review_feature_agreement} compare corresponding PointPCA+ and PointPCA+RS predictors after verifying their intended descriptor order. Across the 160 feature--dataset pairs, median Spearman correlation is 1.000, with a minimum of 0.524. Per-feature errors reveal departures that a single pooled correlation would conceal. On SJTU-PCQA, for example, PointPCA+RS minus PointPCA+ SROCC is -0.006 with a pointwise 95\% interval [-0.015, 0.004]. Table~\ref{tab_review_pairwise} reports paired prediction differences and Holm-adjusted tests; none of the paired comparisons between PointPCA+ and PointPCA+RS is significant after the declared correction. The tests do not establish equivalence; the estimated differences and their intervals quantify the agreement supported by these populations. Inspected source establishes intended feature semantics, while exact historical native binary versions remain a provenance limitation.}
\newcommand{\ReviewIQAInterpretation}{Table~\ref{tab_review_iqa_complementarity} tests whether TOPIQ-FR adds predictive information to DISTS on APSIPA. Adding its six views to the PointPCA$^3$/DISTS representation changes pooled SROCC from 0.937 to 0.928 and RMSE from 0.543 to 0.572. The paired SROCC change is -0.009 (pointwise 95\% interval [-0.020, 0.003], Holm-adjusted $p=1.000$), so the lower point estimate is not a demonstrated performance loss. Combining the two IQAs without PointPCA$^3$ also lowers the SROCC point estimate relative to DISTS alone. High importance in a larger correlated feature set therefore does not imply a benefit from this concatenation. These diagnostic results retain the selected 46-feature design and do not establish behavior beyond APSIPA.}
\newcommand{\ReviewViewInterpretation}{For the fused representation, the median four-view SROCC is 0.934, compared with 0.937 for six views. Across the four pooling functions, fused SROCC ranges from 0.918 to 0.923. Pooling therefore reduces the input dimension from 46 to 41 with lower point estimates in this dataset, while some indexed subsets approach the six-view estimate. Paired reference-clustered intervals against each six-view baseline are retained for every design in the numerical results; these descriptive comparisons do not establish a universally preferable subset or pooling rule.}
\newcommand{\ReviewErrorInterpretation}{Table~\ref{tab_review_error_summary} quantifies dataset-wise errors; WPC2 has the largest mean absolute error relative to its observed MOS range (0.102). In all seven datasets, mean residuals are positive in the lowest MOS tercile and negative in the highest, consistent with compression toward central scores. On WPC2 these biases are 5.488 and -6.242, respectively. The terciles include ties at their empirical boundaries and are descriptive, post-hoc score strata. Tables~\ref{tab_review_cases} and~\ref{tab_review_error_groups} give deterministic examples and the strongest recorded class-level error concentrations. Unknown WPC2 classes remain unknown, and small or constant-score subgroups need not have defined correlation. These observations identify error concentrations without establishing their causes.}
